% --- CORRECCIÓN IMPORTANTE EN LA PRIMERA LÍNEA ---
% Agregamos "xcolor=table" para que funcione \rowcolor sin errores
\documentclass[aspectratio=169, 10pt, xcolor=table]{beamer}

% --- Configuración del Tema (estilo profesional como el ejemplo) ---
\usetheme{Madrid}
\usecolortheme{dolphin}

% --- Paquetes necesarios ---
\usepackage[utf8]{inputenc}
\usepackage[spanish]{babel}
\usepackage{amsmath, amssymb}
\usepackage{booktabs}
\usepackage{graphicx}
\usepackage{listings}
\usepackage{tikz}
\usetikzlibrary{arrows.meta, positioning, shapes.geometric}

% --- Ruta de Imágenes (crea la carpeta "imagenes" y coloca ahí tus archivos) ---
\graphicspath{{./imagenes/}}

% --- Estilo para código Python (configurado para evitar problemas con UTF-8) ---
\definecolor{codegreen}{rgb}{0,0.6,0}
\definecolor{codegray}{rgb}{0.5,0.5,0.5}
\definecolor{codepurple}{rgb}{0.58,0,0.82}
\definecolor{backcolour}{rgb}{0.95,0.95,0.92}
\lstset{
    language=Python,
    backgroundcolor=\color{backcolour},
    commentstyle=\color{codegreen},
    keywordstyle=\color{blue},
    numberstyle=\tiny\color{codegray},
    stringstyle=\color{codepurple},
    basicstyle=\ttfamily\scriptsize,
    breaklines=true,
    frame=single,
    showstringspaces=false,
    literate={á}{{\'a}}1 {é}{{\'e}}1 {í}{{\'i}}1 {ó}{{\'o}}1 {ú}{{\'u}}1 {ñ}{{\~n}}1 {ü}{{\"u}}1
}

% --- Pie de página personalizado ---
\makeatletter

\setbeamertemplate{footline}{
  \leavevmode%
  \hbox{%
  \begin{beamercolorbox}[wd=.333333\paperwidth,ht=2.25ex,dp=1ex,center]{author in head/foot}%
    \usebeamerfont{author in head/foot}\insertshortauthor
  \end{beamercolorbox}%
  \begin{beamercolorbox}[wd=.333333\paperwidth,ht=2.25ex,dp=1ex,center]{title in head/foot}%
    \usebeamerfont{title in head/foot}Sesi\'on 2
  \end{beamercolorbox}%
  \begin{beamercolorbox}[wd=.333333\paperwidth,ht=2.25ex,dp=1ex,right]{date in head/foot}%
    \usebeamerfont{date in head/foot}CÁLCULO Y ÁLGEBRA LINEAL \hfill \insertframenumber/\inserttotalframenumber \hspace*{0.3cm}
  \end{beamercolorbox}}%
  \vskip0pt%
}

\makeatother

% --- Datos de la portada ---
\title[SESI\'ON 2 CÁLCULO]{\textbf{Sesi\'on 2}}
\subtitle{Tensores: Estructuras de Datos Multidimensionales para IA}
\author[Prof. C. S\'anchez]{Carlos C\'esar S\'anchez Coronel}
\institute{\textbf{CÁLCULO Y ÁLGEBRA LINEAL PARA IA}}
\date{}

% Logo opcional (descomentar si tienes la imagen)
% \titlegraphic{\includegraphics[height=1.5cm]{logo.png}}

% --- Eliminar la barra de navegación (opcional) ---
\setbeamertemplate{navigation symbols}{}

% --- Índice al inicio de cada sección ---
\AtBeginSection[]{
  \begin{frame}
    \frametitle{Contenido}
    \tableofcontents[currentsection]
  \end{frame}
}

% ============================================================================
% INICIO DEL DOCUMENTO
% ============================================================================
\begin{document}

% --- Portada ---
\begin{frame}
    \titlepage
\end{frame}

% --- Índice general ---
\begin{frame}{Contenido}
    \tableofcontents
\end{frame}

% ============================================================================
% SECCIÓN 1: ¿QUÉ ES UN TENSOR?
% ============================================================================
\section{¿Qu\'e es un Tensor?}

\begin{frame}{Definici\'on de Tensor}
    \begin{block}{Tensor}
        Un tensor es un arreglo multidimensional de números. Su \textbf{orden} (o rango) es el número de dimensiones.
    \end{block}
    \begin{itemize}
        \item \textbf{Orden 0}: escalar. Ej: $x = 5.2$ (bias de una neurona).
        \item \textbf{Orden 1}: vector. Ej: $\mathbf{x} = [1,2,3]$ (características de una muestra).
        \item \textbf{Orden 2}: matriz. Ej: $\mathbf{X} = \begin{pmatrix}1&2\\3&4\end{pmatrix}$ (imagen en grises 2x2).
        \item \textbf{Orden 3}: imagen RGB: $\mathcal{T} \in \mathbb{R}^{h \times w \times 3}$.
        \item \textbf{Orden 4}: lote de imágenes: $\mathcal{T} \in \mathbb{R}^{b \times h \times w \times c}$.
        \item \textbf{Orden 5}: lote de videos: $\mathcal{T} \in \mathbb{R}^{b \times f \times h \times w \times c}$.
    \end{itemize}
\end{frame}

\begin{frame}{Representaci\'on visual de tensores}
    \begin{center}
        \begin{tikzpicture}[scale=0.7]
            % Escalar
            \draw[fill=blue!20] (0,1) circle (0.5);
            \node at (0,1) {$5.2$};
            \node at (0,2) {Escalar (0D)};
            % Vector
            \draw[fill=green!20] (3,0) rectangle (5,2);
            \node at (4,1) {$[1,2,3]$};
            \node at (4,2.5) {Vector (1D)};
            % Matriz
            \draw[fill=red!20] (7,0) rectangle (10,2);
            \draw (7.5,0.5) -- (7.5,2);
            \draw (8.5,0.5) -- (8.5,2);
            \draw (9.5,0.5) -- (9.5,2);
            \node at (8,1.2) {$a$};
            \node at (9,1.2) {$b$};
            \node at (8,0.7) {$c$};
            \node at (9,0.7) {$d$};
            \node at (8.5,2.5) {Matriz (2D)};
            % Tensor 3D
            \draw[fill=yellow!20] (0,-3) rectangle (2,0);
            \draw (0.5,-3) -- (0.5,0);
            \draw (1,-3) -- (1,0);
            \draw (1.5,-3) -- (1.5,0);
            \node at (1,-1.5) {Tensor 3D};
        \end{tikzpicture}
    \end{center}
\end{frame}

% ============================================================================
% SECCIÓN 2: NOTACIÓN EN PAPERS DE IA
% ============================================================================
\section{Notaci\'on en Papers de IA}

\begin{frame}{Interpretando la notación de tensores}
    \begin{itemize}
        \item $\mathbf{X} \in \mathbb{R}^{n \times d}$: matriz con $n$ muestras y $d$ características.
        \item $\mathcal{T} \in \mathbb{R}^{b \times c \times h \times w}$: tensor de imágenes (batch, canales, alto, ancho).
        \item $x_i$: $i$-ésimo elemento de un vector.
        \item $X_{ij}$: elemento en fila $i$, columna $j$ de una matriz.
        \item $\mathcal{T}_{ijk}$ o $\mathcal{T}_{ij,k}$: elemento de un tensor 3D.
    \end{itemize}
\end{frame}

\begin{frame}{Ejemplo de notación}
    \begin{exampleblock}{Imágenes}
        $\mathbf{X} \in \mathbb{R}^{32 \times 3 \times 224 \times 224}$ representa un lote de 32 imágenes RGB de 224x224 píxeles.
        \begin{itemize}
            \item 32 = batch size
            \item 3 = canales (R,G,B)
            \item 224 = alto
            \item 224 = ancho
        \end{itemize}
    \end{exampleblock}
    \begin{exampleblock}{Texto}
        $\mathbf{E} \in \mathbb{R}^{10 \times 300}$: una frase de 10 palabras con embeddings de dimensión 300.
    \end{exampleblock}
\end{frame}

% ============================================================================
% SECCIÓN 3: OPERACIONES FUNDAMENTALES
% ============================================================================
\section{Operaciones Fundamentales con Tensores}

\subsection{Transposición}

\begin{frame}{Transposición de matrices y tensores}
    \begin{block}{Matriz}
        $(A^T)_{ij} = A_{ji}$
        \[
        A = \begin{pmatrix}1&2\\3&4\\5&6\end{pmatrix} \quad A^T = \begin{pmatrix}1&3&5\\2&4&6\end{pmatrix}
        \]
    \end{block}
    \begin{block}{Tensor 3D}
        Podemos permutar ejes. Ej: tensor de forma $(2,3,4)$ a $(4,2,3)$ con \texttt{tensor.transpose(2,0,1)}.
    \end{block}
\end{frame}

\subsection{Redimensionamiento (Reshape)}

\begin{frame}{Cambiar la forma sin alterar los datos}
    \begin{itemize}
        \item Ejemplo: imagen de $28 \times 28$ se aplana a un vector de $784$ elementos.
        \item Condición: el número total de elementos debe conservarse.
        \item En NumPy: \texttt{vector.reshape(2,3)} si el vector tiene 6 elementos.
    \end{itemize}
    \begin{exampleblock}{Ejemplo}
        \[
        \mathbf{v} = [1,2,3,4,5,6] \quad \rightarrow \quad \begin{pmatrix}1&2&3\\4&5&6\end{pmatrix}
        \]
    \end{exampleblock}
\end{frame}

\subsection{Normas de Vectores}

\begin{frame}{Normas L1, L2 e infinito}
    \begin{columns}[t]
        \column{0.33\textwidth}
        \textbf{L2 (Euclidiana)}
        \[
        \|\mathbf{x}\|_2 = \sqrt{\sum_i x_i^2}
        \]
        \column{0.33\textwidth}
        \textbf{L1 (Manhattan)}
        \[
        \|\mathbf{x}\|_1 = \sum_i |x_i|
        \]
        \column{0.33\textwidth}
        \textbf{L-infinito}
        \[
        \|\mathbf{x}\|_\infty = \max_i |x_i|
        \]
    \end{columns}
    \vspace{0.3cm}
    \begin{exampleblock}{Ejemplo con $\mathbf{x} = (3, -4)$}
        \begin{align*}
            \|x\|_2 &= \sqrt{3^2+(-4)^2}=5 \\
            \|x\|_1 &= 3+4=7 \\
            \|x\|_\infty &= \max(3,4)=4
        \end{align*}
    \end{exampleblock}
\end{frame}

\begin{frame}{Interpretación geométrica de las normas}
    \begin{center}
        \begin{tikzpicture}[scale=0.8]
            % Círculo L2
            \draw[thick, blue] (0,0) circle (2);
            \node at (0,0.5) {$L_2$};
            % Diamante L1
            \draw[thick, red] (4,0) -- (6,2) -- (8,0) -- (6,-2) -- cycle;
            \node at (6,0.5) {$L_1$};
            % Cuadrado L-inf
            \draw[thick, green] (10,2) -- (14,2) -- (14,-2) -- (10,-2) -- cycle;
            \node at (12,0.5) {$L_\infty$};
        \end{tikzpicture}
    \end{center}
    \begin{itemize}
        \item Conjuntos de vectores con norma 1.
        \item L1 produce soluciones dispersas.
        \item L2 es el círculo (rotacionalmente invariante).
        \item L-inf es el cuadrado (ignora valores pequeños).
    \end{itemize}
\end{frame}

% ============================================================================
% SECCIÓN 4: APLICACIONES EN IA
% ============================================================================
\section{Aplicaciones en IA}

\subsection{Visión por Computador}

\begin{frame}{Tensores en imágenes y videos}
    \begin{itemize}
        \item Imagen en escala de grises: matriz $\in \mathbb{R}^{h \times w}$.
        \item Imagen RGB: tensor $\in \mathbb{R}^{h \times w \times 3}$.
        \item Lote de imágenes: tensor $\in \mathbb{R}^{b \times h \times w \times c}$.
        \item Lote de videos: $\in \mathbb{R}^{b \times f \times h \times w \times c}$.
    \end{itemize}
\end{frame}

\subsection{Procesamiento de Lenguaje Natural}

\begin{frame}{Tensores en NLP}
    \begin{itemize}
        \item Embedding de palabra: vector $\in \mathbb{R}^{d}$.
        \item Frase (secuencia): matriz $\in \mathbb{R}^{n \times d}$.
        \item Lote de frases: tensor $\in \mathbb{R}^{b \times n \times d}$.
        \item Atención en Transformers: tensores $(batch, heads, seq, seq)$.
    \end{itemize}
\end{frame}

\subsection{Regularización en Machine Learning}

\begin{frame}{Regularización L1 y L2}
    \begin{block}{Ridge (L2)}
        \[
        \mathcal{L} = \text{MSE} + \lambda \|\mathbf{w}\|_2^2
        \]
        \begin{itemize}
            \item Favorece pesos pequeños pero no nulos.
            \item Reduce la varianza del modelo.
        \end{itemize}
    \end{block}
    \begin{block}{Lasso (L1)}
        \[
        \mathcal{L} = \text{MSE} + \lambda \|\mathbf{w}\|_1
        \]
        \begin{itemize}
            \item Puede llevar pesos exactamente a cero.
            \item Realiza selección automática de características.
        \end{itemize}
    \end{block}
\end{frame}

% ============================================================================
% SECCIÓN 5: IMPLEMENTACIÓN EN PYTHON CON NUMPY
% ============================================================================
\section{Python con NumPy: Manipulación de Tensores}

\begin{frame}[fragile]{Creación de tensores}
    \begin{lstlisting}
import numpy as np

# Escalar (0D)
escalar = np.array(5)
print(f"Escalar: {escalar}, dimensiones: {escalar.ndim}")

# Vector (1D)
vector = np.array([1,2,3,4])
print(f"Vector shape: {vector.shape}")

# Matriz (2D)
matriz = np.array([[1,2],[3,4],[5,6]])
print(f"Matriz shape: {matriz.shape}")

# Tensor 3D (2 imagenes 3x3)
tensor_3d = np.random.rand(2,3,3)
print(f"Tensor 3D shape: {tensor_3d.shape}")

# Tensor 4D (lote de 2 imagenes RGB 32x32)
tensor_4d = np.random.rand(2,32,32,3)
print(f"Tensor 4D shape: {tensor_4d.shape}")
    \end{lstlisting}
\end{frame}

\begin{frame}[fragile]{Atributos importantes}
    \begin{lstlisting}
a = np.array([[1,2,3],[4,5,6]])
print(f"Forma: {a.shape}")        # (2,3)
print(f"Dimensiones: {a.ndim}")    # 2
print(f"Tamaño total: {a.size}")   # 6
print(f"Tipo de datos: {a.dtype}") # int64
    \end{lstlisting}
\end{frame}

\begin{frame}[fragile]{Transposición y reshape}
    \begin{lstlisting}
# Transposición
matriz = np.array([[1,2],[3,4],[5,6]])
print("Transpuesta:\n", matriz.T)

# Permutar ejes en tensor 3D
tensor = np.random.rand(2,3,4)
permutado = tensor.transpose(2,0,1)
print(f"Original shape: {tensor.shape}")
print(f"Nuevo shape: {permutado.shape}")  # (4,2,3)

# Redimensionamiento (reshape)
vector = np.array([1,2,3,4,5,6])
matriz_2x3 = vector.reshape(2,3)
print(matriz_2x3)
# Usar -1 para inferir dimensión
otra = vector.reshape(3,-1)  # 3 filas, 2 columnas
print(otra)
    \end{lstlisting}
\end{frame}

\begin{frame}[fragile]{Cálculo de normas}
    \begin{lstlisting}
v = np.array([3, -4])

norma_l2 = np.linalg.norm(v, ord=2)    # 5.0
norma_l1 = np.linalg.norm(v, ord=1)    # 7.0
norma_inf = np.linalg.norm(v, ord=np.inf) # 4.0
norma_l2_cuad = np.sum(v**2)            # 25

print(f"L2: {norma_l2}, L1: {norma_l1}, L-inf: {norma_inf}")

# Normas por filas en una matriz
matriz = np.array([[1,2],[3,4],[5,6]])
normas_filas = np.linalg.norm(matriz, axis=1)
print("Normas por fila:", normas_filas)
    \end{lstlisting}
\end{frame}

% ============================================================================
% SECCIÓN 6: EJEMPLO INTEGRADO
% ============================================================================
\section{Ejemplo Integrado: Procesamiento de un Lote de Imágenes}

\begin{frame}[fragile]{Simulación de lote de imágenes}
    \begin{lstlisting}
import numpy as np

# Crear lote de 3 imagenes en grises de 4x4 (valores 0-255)
batch = np.random.randint(0, 256, size=(3,4,4)).astype(np.float32)
print(f"Forma del lote: {batch.shape}")  # (3,4,4)

# Aplanar cada imagen para calcular norma
imagenes_aplanadas = batch.reshape(3, -1)
normas = np.linalg.norm(imagenes_aplanadas, axis=1)
print("Normas L2 de cada imagen:", normas)

# Normalizar cada imagen (dividir por su norma)
batch_normalizado = batch / normas[:, np.newaxis, np.newaxis]

# Verificar normas después de normalizar (deben ser 1)
nuevas_normas = np.linalg.norm(batch_normalizado.reshape(3,-1), axis=1)
print("Nuevas normas (deberian ser 1):", nuevas_normas)
    \end{lstlisting}
\end{frame}

% ============================================================================
% SECCIÓN 7: INTERPRETANDO NOTACIÓN EN PAPERS
% ============================================================================
\section{Interpretando Notación en Papers}

\begin{frame}{Ejemplo de ecuación con tensores}
    \begin{block}{Atención en Transformers}
        \[
        \text{Attention}(Q, K, V) = \text{softmax}\left(\frac{QK^\top}{\sqrt{d_k}}\right)V
        \]
    \end{block}
    \begin{itemize}
        \item $Q, K, V$ son tensores de dimensión $(batch, seqlen, d_{model})$.
        \item $QK^\top$ produce un tensor de atención $(batch, seqlen, seqlen)$.
        \item El producto final combina los valores ponderados.
    \end{itemize}
\end{frame}

\begin{frame}{Ejemplo de función de pérdida regularizada}
    \[
    \mathcal{L}(\mathbf{w}) = \frac{1}{N}\sum_{i=1}^{N} \|\mathbf{y}_i - \hat{\mathbf{y}}_i\|_2^2 + \lambda \|\mathbf{w}\|_1
    \]
    \begin{itemize}
        \item $\mathbf{y}_i, \hat{\mathbf{y}}_i$ son vectores de etiquetas y predicciones.
        \item $\|\cdot\|_2^2$ es la norma L2 al cuadrado (suma de cuadrados).
        \item $\|\mathbf{w}\|_1$ es la norma L1 (regularización Lasso).
        \item $\lambda$ controla la fuerza de la regularización.
    \end{itemize}
\end{frame}

% ============================================================================
% SECCIÓN 8: RESUMEN
% ============================================================================
\section{Resumen}

\begin{frame}{Lo que hemos aprendido}
    \begin{exampleblock}{Conceptos clave}
        \begin{itemize}
            \item \textbf{Tensores}: generalización de escalares, vectores y matrices a más dimensiones.
            \item \textbf{Notación}: $\mathbb{R}^{d_1\times\cdots\times d_k}$ y cómo interpretarla en papers.
            \item \textbf{Operaciones}: transposición, reshape, normas L1, L2, L∞.
            \item \textbf{Aplicaciones}: visión (imágenes), NLP (embeddings), regularización (Ridge, Lasso).
            \item \textbf{NumPy}: creación, atributos, transposición, reshape, cálculo de normas.
            \item \textbf{Caso integrado}: procesamiento y normalización de un lote de imágenes.
            \item \textbf{Lectura de ecuaciones}: atención, pérdida regularizada.
        \end{itemize}
    \end{exampleblock}
\end{frame}

\begin{frame}{Conexión con cursos futuros}
    \begin{itemize}
        \item Estos conceptos son la base para:
        \begin{itemize}
            \item Redes convolucionales (CNN): tensores de imágenes.
            \item Transformers: tensores de atención.
            \item Regularización en modelos lineales y redes profundas.
            \item Preprocesamiento de datos en pipelines de ML.
        \end{itemize}
    \end{itemize}
\end{frame}

% --- Diapositiva final ---
\begin{frame}
    \centering
    \Huge \textbf{¡Gracias!}

\end{frame}

\end{document}